\begin{table}[h]
    \centering
    \caption{EU ETS coverage of Belgian CO\textsubscript{2} emissions in our sample, selected years}
    \label{tab:ets_coverage_belgium}
    \small
    \begin{tabular}{l c c c c c}
        \toprule
         & \multicolumn{3}{c}{Emissions (Mt CO\textsubscript{2})} & \multicolumn{2}{c}{Coverage (\%)} \\
        \cmidrule(lr){2-4} \cmidrule(lr){5-6}
        Year & Sample EU ETS & Belgian total & Stationary & of total & of stationary \\
        \midrule
        2005 & 50.5 & 125.7 & 77.1 & 40 & 65 \\
        2008 & 50.4 & 120.2 & 71.7 & 42 & 70 \\
        2012 & 38.7 & 102.5 & 58.9 & 38 & 66 \\
        2020 & 39.4 & 91.3 & 54.0 & 43 & 73 \\
        \bottomrule
    \end{tabular}
    \begin{minipage}{\textwidth}
        {\footnotesize \vspace{0.25cm} \noindent \textit{Notes:} \emph{Sample EU ETS} sums verified CO\textsubscript{2} emissions of EU ETS firms in our analysis sample (firms meeting the De Loecker/Dhyne sample-selection criteria) from installations located in Belgium. \emph{Belgian total} is Belgium's total CO\textsubscript{2} emissions excluding land use, land-use change and forestry, from the National Inventory Report (NIR). \emph{Stationary} sums CO\textsubscript{2} emissions from stationary installations as reported in the NIR --- fuel combustion in energy industries (Common Reporting Format category 1.A.1), manufacturing and construction (1.A.2), other sectors excluding residential heating (1.A.4 -- 1.A.4.b), plus all industrial process emissions (2). Coverage shares are sample EU ETS emissions divided by each total.}
    \end{minipage}
\end{table}
